% ============================================================================
\section{Implementation and Supporting Derivations}
% ============================================================================

\subsection{Configuration}
\label{app:reference_tables}

\begin{table}[H]
\centering
\footnotesize
\caption{Configuration of the four-domain study. Entries marked TODO
require the corresponding training or evaluation artifact.}
\label{tab:pipeline}
\begin{tabular}{p{0.19\linewidth}p{0.73\linewidth}}
\toprule
Component & Configuration \\
\midrule
Student & Qwen3-1.7B \citep{qwen2025qwen3}, non-thinking mode;
\missing{frozen model and tokenizer revisions} \\
Teachers & Domain-specific Qwen3-1.7B RL checkpoints for math and
instruction following; one shared Qwen3-4B weight set for code and science.
All use the same vocabulary and tokenizer.
\missing{exact checkpoint identifiers, RL training recipes, and scores} \\
Training data & $16{,}000$ unique prompts per task, shuffled once and
consumed without replacement.
\missing{datasets, versions, licenses, filters, and held-out split} \\
Default loss & Teacher top-$64$ corrected reverse KL
(Equation~\ref{eq:dense_topk_loss}); full-vocabulary probabilities;
position mean within response, response mean within task; $w_i=1/4$ \\
Rollout & One response per prompt and one update per fresh batch;
maximum response length $4{,}096$; truncated responses retained.
\missing{temperature, EOS handling, and decoding settings} \\
Step and budget & $b=4$, $G=16$, $2\leq m_i\leq8$, $500$ steps,
$32{,}000$ attempted responses per run, one seed per method \\
Optimizer & AdamW;
\missing{learning rate and schedule, betas, epsilon, weight decay,
clipping, precision, and memory configuration} \\
GPAS & Pre-step bias-corrected Adam second moment for $D$;
noise and time EMA decay $0.9$; first step Uniform;
bounded proportional rounding for per-step counts and integer
enumeration for cost-aware counts \\
Evaluation & $128$ held-out prompts per task every $50$ steps;
MATH-500, LiveCodeBench, IFBench, and GPQA-Diamond.
\missing{versions, evaluation seeds, and decoding protocols} \\
System & One 96\,GB learner GPU and one 48\,GB rollout/teacher GPU;
the code and science endpoints share weights.
\missing{GPU models, inference and training software revisions,
peak memory, and GPU hours} \\
\bottomrule
\end{tabular}
\end{table}

\subsection{Dense loss implementation and online statistics}
\label{app:implementation}

At each student-generated prefix, the teacher returns the identifiers
and full-vocabulary log probabilities of its top-$64$ tokens. The
student computes its full-vocabulary log normalizer and gathers its
probabilities at those identifiers. Autodifferentiation is applied to
the whole expression in Equation~\ref{eq:dense_topk_loss}, including
the $-p_\theta$ term. The retained probabilities are not renormalized
within the top-$k$ set, and the loss does not use a sampled tail
correction. The operation is dense over the retained set and uses
sparse teacher output.

GPAS collects statistics from unweighted micro-batch gradients. For a
task block, initialize its running mean $h=0$ and scalar $S=0$.
After micro-batch $s$, form $\delta=g_{i,s}-h$, then update
\begin{equation}
    h\leftarrow h+\delta/s,\qquad
    S\leftarrow S+\frac{s-1}{s}\|D\delta\|_2^2.
    \label{eq:welford}
\end{equation}
At the end of the block, $\widehat e_i=S/(m_i-1)$ and
$A\leftarrow A+w_i h$. The mean buffer is reused for the next task.
The same recurrence with $D=I$ records unpreconditioned noise.
Reductions use float32 accumulation. Gradients are unscaled before
measurement if mixed-precision loss scaling is active, and clipping
is applied to the assembled gradient $A$.

The optimizer state is fixed during this collection. Before an Adam
second moment is available, the first Uniform step uses $D=I$.
Subsequent steps use the pre-step bias-corrected second moment.
The first observation initializes each noise EMA; later observations
use decay $0.9$. The scheduler reads the completed step's averages
before generating the next step's data.

For per-step allocation, set $a_i=w_i\sqrt{\widetilde e_i}$ and find
a scale $\lambda$ such that
$\sum_i\operatorname{clip}(\lambda a_i,m_{\min},m_{\max})=G$.
Floor these counts and assign remaining units in descending order
of fractional remainder, respecting the bounds; ties follow the
fixed task order. Zero-noise tasks receive the lower bound unless
the other tasks reach their upper bounds; if all scores are zero,
use Uniform. The cost-aware scheduler evaluates all feasible integer
vectors using the smoothed noise and timing estimates.

\subsection{Comparison recipes}
\label{app:baselines}

\paragraph{D$^3$-MOPD scheduling.}
We use the remaining-gap and descent-velocity construction of
\citet{sun2026d3mopd}. Let $\bar L_i(t)$ be the smoothed task loss
and $L_i^0$ the mean of its first five observations. With $W=10$
and up to $R=3$ available windows, the composite signal is
\begin{equation}
\begin{split}
    r_i(t)&=\bar L_i(t)/L_i^0,\\
    v_i(t)&=\max\left\{0,\frac1{R'}
       \sum_{j=0}^{R'-1}
       \frac{\bar L_i(t-(j+1)W)-\bar L_i(t-jW)}
            {\bar L_i(t-(j+1)W)}\right\},\\
    s_i(t)&=r_i(t)v_i(t).
\end{split}
\end{equation}
We retain the published max-normalization, softmax temperature $0.5$,
per-domain probability floor $0.10$, and multiplicative batch jitter
of amplitude $0.30$. The scheduler refreshes every $10$ steps,
using the paper's history warmup and EMA window of $10$.
All KL denominators use the published floor of $0.15$.
The signal is computed from the default dense task loss in our pipeline.
The resulting mixture is converted to micro-batch counts with the
shared bounds $[2,8]$. Thus this is a controlled adaptation to our loss
scale, micro-batch unit, and budget.

The two D$^3$ rows use the same scheduling rule. The fixed-weight
row assembles $A=\sum_iw_i\bar g_i$; the recipe-level row assembles
$A=\sum_i(m_i/G)\bar g_i$. This makes the effect of task weighting
visible alongside the comparison of scheduling signals.

\paragraph{Open-MOPD with one update per rollout.}
We follow the released student top-$16$ dense candidate loss,
token-share balancing, and gap-aware weighting
\citep{gao2026openmopddiagnosingfixingcapability,openmopd2026code}.
The released implementation multiplies candidate-wise detached
advantages by candidate-wise student probability ratios and then
sums over candidates. Each retained token therefore contributes its own gradient. In our setting, the initial task target is uniform, and each
batch contains four micro-batches per task. The recipe's token and
gap weights remain active. One optimizer step per rollout batch
corresponds to $K=1$, so there is no reuse-induced reward staleness.
\missing{Record the exact configuration for gap smoothing, floors,
and weight normalization from the integrated implementation.}

\paragraph{Tail-aware supervision.}
TA-OPD \citep{huang2026tailaware} groups the omitted vocabulary
into one tail event. For teacher support $S_i(c)$, write
$p_{\rm tail}=1-\sum_{a\in S_i(c)}p_\theta(a\mid c)$ and
$q_{\rm tail}=1-\sum_{a\in S_i(c)}q_i(a\mid c)$. Its position loss is
\begin{equation}
    \ell_i^{\rm TA}=
       \sum_{a\in S_i(c)}p_\theta(a\mid c)
                     \log\frac{p_\theta(a\mid c)}{q_i(a\mid c)}
       +p_{\rm tail}\log\frac{p_{\rm tail}}{q_{\rm tail}}.
\end{equation}
Both runs use $k=64$, the same response averaging and fixed task
weights, and the same training budget as the core comparison.
Only the allocation differs between TA-OPD + Uniform and
TA-OPD + GPAS.

\subsection{Step-local objective}
\label{app:local_objective}

Let $\theta_t$ be the current student and $\mathcal R_i^t$ its
distribution over task-$i$ prompts and responses at step $t$.
The local objective differentiated by dense OPD is
\begin{equation}
    F_t(\theta)=
       \sum_i w_i\,\mathbb E_{(x,y)\sim\mathcal R_i^t}
       \left[\frac1{|y|}\sum_{u=1}^{|y|}
            \ell_i^{(k)}(\theta;x,y_{<u})\right].
\end{equation}
The rollout distribution is held fixed in this derivative. Thus the
mean micro-batch update at $\theta_t$ is $\nabla F_t(\theta_t)$.
The distribution is refreshed at the next step. Fixing $w_i$ preserves
the chosen task weighting through this sequence of local objectives.

\subsection{Variance and the allocation rule}
\label{app:allocation_proofs}

Condition on the current state and pre-step scaling $D$. In the
independent micro-batch model, cross-task centered terms have zero
expectation, giving
\begin{equation}
    \mathbb E\|D(A-\mu)\|_2^2
       =\sum_i w_i^2\mathbb E\|D(\bar g_i-\mu_i)\|_2^2
       =\sum_i\frac{w_i^2e_i}{m_i}.
\end{equation}
Equation~\ref{eq:noise_estimate} is the ordinary sample variance
summed over preconditioned coordinates.
Writing $a_i=w_i\sqrt{e_i}$, Cauchy--Schwarz gives
\begin{equation}
    \left(\sum_i m_i\right)
    \left(\sum_i\frac{a_i^2}{m_i}\right)
    \geq\left(\sum_i a_i\right)^2.
\end{equation}
Equality holds at $m_i\propto a_i$, yielding the continuous
per-step rule. With unequal costs and zero fixed step time,
\begin{equation}
    \left(\sum_i m_i\tau_i\right)
    \left(\sum_i\frac{a_i^2}{m_i}\right)
    \geq\left(\sum_i a_i\sqrt{\tau_i}\right)^2,
\end{equation}
with equality at $m_i\propto a_i/\sqrt{\tau_i}$.
Positive fixed time and integer bounds are handled by the finite
search used in the cost-aware implementation.

\subsection{Connection to a local update}
\label{app:descent}

For an $L$-smooth $F_t$ and the frozen-preconditioner step
$\theta'=\theta-\eta DA$, let $\mu=\nabla F_t(\theta)$.
Smoothness and the bias--variance identity give
\begin{equation}
    \mathbb E F_t(\theta')
       \leq F_t(\theta)-\eta\mu^\top D\mu
       +\frac{\eta^2L}{2}
         \left(\|D\mu\|_2^2+V(m)\right).
\end{equation}
At the fixed state, the allocation-dependent term is $V(m)$.
This motivates measuring noise in update coordinates. GPAS uses
this local criterion with ordinary AdamW; momentum, clipping, and
the new second-moment update remain part of that optimizer.

\subsection{Timing and released traces}
\label{app:system_details}

Task blocks record rollout, teacher-scoring, backward, and statistics
time, together with response lengths and peak memory. The marginal
per-micro-batch estimates form $\tau_i$; synchronization, optimizer
work, and other count-independent step operations form $C$.
Their EMA estimates use decay $0.9$. We compare predicted and
observed step times and evaluate efficiency with actual elapsed
time on both reserved GPUs. Evaluation compute is reported separately.
If a deployment overlaps task blocks, its timing predictor should
represent that execution schedule.

Checkpoints retain optimizer state, task stream positions, counts,
and noise and time averages. The released traces include per-task
held-out loss, prompt consumption, response length and truncation,
micro-batch counts, noise estimates, and measured step time.
\missing{Add the artifact location and populate the configuration
and empirical figures from these records.}
